%% BioAgent preprint -- NeurIPS/ICLR-style preprint format
%%
%% Compile:
%%   pdflatex manuscript.tex
%%   bibtex   manuscript
%%   pdflatex manuscript.tex && pdflatex manuscript.tex

\documentclass[11pt]{article}

%% ── Packages ─────────────────────────────────────────────────────────────────
\usepackage[utf8]{inputenc}
\usepackage[T1]{fontenc}
\usepackage{mathptmx}          % Times Roman text + matching math
\usepackage[scaled=0.92]{helvet}
\usepackage{microtype}
\usepackage{graphicx}
\usepackage{amsmath,amssymb}
\usepackage{booktabs}
\usepackage[hyphens]{url}
\usepackage{hyperref}
\usepackage{xcolor}
\usepackage{listings}
\usepackage[numbers,square,sort&compress]{natbib}
\usepackage{geometry}
\usepackage{titlesec}
\usepackage{titling}
\usepackage{caption}

%% NeurIPS-like single-column page geometry
\geometry{
  letterpaper,
  left=1.25in, right=1.25in,
  top=1in, bottom=1in,
}
\setlength{\emergencystretch}{3em}
\sloppy

\captionsetup{font=small, labelfont=bf, labelsep=period}

%% Section heading style: sans-serif bold, flush-left
\titleformat{\section}{\normalfont\Large\bfseries\sffamily}{\thesection}{1em}{}
\titleformat{\subsection}{\normalfont\large\bfseries\sffamily}{\thesubsection}{1em}{}
\titleformat{\subsubsection}{\normalfont\normalsize\bfseries\sffamily}{\thesubsubsection}{1em}{}
\titleformat{\paragraph}[runin]{\normalfont\normalsize\bfseries}{}{0pt}{}[.]
\titlespacing*{\section}{0pt}{2.0ex plus 0.5ex}{1.0ex plus 0.2ex}
\titlespacing*{\subsection}{0pt}{1.5ex plus 0.4ex}{0.8ex plus 0.2ex}

\hypersetup{
  colorlinks=true,
  linkcolor=black,
  citecolor=blue!55!black,
  urlcolor=blue!55!black,
}

%% NeurIPS-style abstract (centered heading, indented narrower body)
\renewenvironment{abstract}{%
  \begin{center}{\bfseries\large Abstract}\end{center}
  \begin{list}{}{%
    \setlength{\leftmargin}{0.4in}%
    \setlength{\rightmargin}{0.4in}}\item\relax\small
}{%
  \end{list}\vspace{0.5ex}%
}

%% ── Title / Authors (NeurIPS-style centered block) ───────────────────────────
\pretitle{\begin{center}\LARGE\bfseries}
\posttitle{\par\end{center}\vskip 0.5em}
\preauthor{\begin{center}\large\lineskip 0.5em\begin{tabular}[t]{c}}
\postauthor{\end{tabular}\par\end{center}}
\predate{}\postdate{}

\title{BioAgent: An Autonomous Multi-Agent System for\\
  End-to-End Bioinformatics Research}

\author{%
  Nigmat Rahim\thanks{Corresponding author.
    ORCID: \href{https://orcid.org/0009-0009-9886-5224}{0009-0009-9886-5224}.} \\
  School of Life Sciences \\
  Peking University \\
  Beijing 100871, China \\
  \texttt{nigmatrahim@stu.pku.edu.cn} \\
}

\date{}

\begin{document}

\maketitle

%% ── Abstract ─────────────────────────────────────────────────────────────────
\begin{abstract}
\textbf{Motivation:}
The pace of biomedical discovery is increasingly bottlenecked by the manual
effort required to synthesise literature, form testable hypotheses, execute
computational analyses and produce publication-quality manuscripts.
Large language models (LLMs) enable new automation opportunities, but
end-to-end autonomous research remains difficult: existing systems either
address isolated stages or operate on pre-curated data and therefore never
encounter the brittleness of real biomedical data portals (partial
downloads, corrupted archives, and geography-dependent latency).

\textbf{Results:}
We present \textbf{BioAgent}, an autonomous multi-agent system that
orchestrates seven specialised agents (Orchestrator, Literature, Planner,
DataAcquisition, Analyst, Writer, Visualization, and Review) in a directed
LangGraph workflow to conduct end-to-end bioinformatics research.
Given a natural-language research question, BioAgent autonomously retrieves
and synthesises primary literature via BioMCP (PubMed, ClinVar, OncoKB),
generates novel testable hypotheses, downloads real datasets from GEO,
cBioPortal, GDC, NCBI and ENCODE with per-file mirror provenance, executes
reproducible statistical analyses (seeded randomness, SHA-256 hashing),
drafts full IMRAD manuscripts with PMID-grounded citations, generates
publication-quality figures, and self-reviews the output against a
six-dimension rubric.
To our knowledge, BioAgent is the first autonomous biomedical research
agent that closes the loop from \emph{live} multi-repository data
acquisition (with mirror-aware, retry/resume-hardened downloads and
per-file \texttt{data\_source} provenance) through analysis to a
self-reviewed IMRAD manuscript.
Comparable recent systems (Biomni, CellAgent, GeneAgent, TxAgent)
either operate on pre-curated or user-supplied data and stop at
analysis output.
A distinctive contribution is the resilient data-acquisition layer
(\texttt{httpx}+\texttt{tenacity} with Range-resume, gzip-integrity
validation, adaptive timeouts, and mirror-first routing through EBI
ArrayExpress, ENA direct \texttt{fastq.gz}, and the 10x~Genomics
Cloudflare CDN) that converts geography-hostile endpoints from a dominant
failure mode into transparent fallback.
Across three benchmark cases (BRAF V600E melanoma, TP53 pan-cancer, and
10x~Genomics PBMC~3k), weighted scores on a 0--10 rubric improved from
a previous baseline of 4.27 to 7.65 on average, with the TP53 case rising
from a pre-fix 1.06 (aborted in literature phase) to 8.42 with a
4{,}439-word manuscript and six figures.
Full pipelines complete in 45--82 minutes at \$7.62--\$14.19 per case,
orders of magnitude below equivalent expert-researcher time cost, with
every artefact traceable through a serialised \texttt{ResearchState}
blackboard.

\textbf{Availability:}
BioAgent is open-source (MIT) at
\url{https://github.com/Nigmat-future/lyceumai}.
A Docker image, pinned dependency lock file, and one-command reproduction
script (\texttt{scripts/reproduce\_benchmark.sh}) are provided.
Requires Python~$\geq$3.11 and an Anthropic-compatible LLM endpoint.
\end{abstract}

\noindent
\textbf{Keywords:} autonomous agents, large language models, bioinformatics
automation, LangGraph, multi-agent systems, reproducible research

%% ─────────────────────────────────────────────────────────────────────────────
\section{Introduction}

The volume of biomedical literature doubles approximately every nine
years~\cite{bornmann2015}, making comprehensive literature review an
increasingly intractable manual task.
At the same time, reproducibility in computational biology remains a
challenge~\cite{Piccolo2016,Baker2016}: analyses are often
underdocumented, tools version-pinned imprecisely, and random seeds
omitted.

Recent advances in large language models~\cite{OpenAI2023gpt4,anthropic2024}
have spurred development of AI-assisted research tools, including
biomedical-domain variants~\cite{Lee2023biomedlm}.
However, existing solutions, including web-search augmented chatbots and
single-prompt code generators, are stateless, lack citation grounding, and
require continuous human oversight.
What is missing is a \emph{full-pipeline} system that maintains shared
research state across agents, integrates authoritative biomedical databases,
produces reproducible code, and outputs artefacts that can be submitted to
a journal with minimal human editing.

BioAgent addresses this gap.
Its key contributions are:

\begin{enumerate}
  \item A seven-agent LangGraph \texttt{StateGraph} with blackboard
        architecture (\texttt{ResearchState}, 30+ typed fields) and
        SQLite-backed session checkpointing that supports mid-run
        resumption after API-gateway or network failures
        (\texttt{benchmarks/resume\_run.py}; demonstrated to save
        $\sim$40~min and $\sim$\$7 per recovered case).
  \item A dedicated \textbf{DataAcquisitionAgent} spanning nine
        primary sources (GEO, cBioPortal, GDC/TCGA, NCBI E-utilities,
        ENCODE, ENA, EBI ArrayExpress, 10x~Genomics CDN, generic URL)
        with a resilient HTTP backbone
        (\texttt{bioagent/tools/data/\_http.py}) featuring
        exponential-backoff retry, HTTP-Range-based resume, adaptive
        read timeouts derived from \texttt{Content-Length}, and
        post-download gzip-integrity validation. To our knowledge this
        is the first autonomous research agent treating mirror choice
        as first-class provenance (\S\ref{sec:resilience}), recording
        which mirror served each file in a \texttt{data\_source} field
        that the WriterAgent cites in Methods.
  \item An \textbf{orchestrator loop detector} (\S\ref{sec:orchestrator})
        combining a deterministic fallback (forced forward progression
        after three consecutive same-phase selections) with a
        phase-history-aware prompt rule, which eliminated a regression
        in which the LLM re-selected
        \texttt{hypothesis\_generation} indefinitely after
        \texttt{code\_execution}.
  \item BioMCP integration for structured literature retrieval with
        PMID-level citation tracking; an export pipeline producing
        \LaTeX{} manuscripts in Bioinformatics~(OUP) format with
        auto-generated BibTeX; and reproducible Python execution with
        automatic seed injection and SHA-256-hashed provenance.
  \item A six-dimension evaluation framework (literature coverage,
        hypothesis quality, analysis correctness, writing quality,
        figure quality, efficiency) validated on three benchmark cases
        with before/after comparisons against the pre-resilience
        pipeline (v0.1), showing mean improvement from 4.27 to 7.65 on
        a 0--10 rubric; the TP53 case rose from 1.06 (no manuscript)
        to 8.42 (full IMRAD + 6 figures + passing self-review), the
        strongest single-variable quantitative argument for treating
        data-acquisition resilience as a first-class component of
        autonomous research agents rather than as plumbing.
\end{enumerate}

%% ─────────────────────────────────────────────────────────────────────────────
\section{System Architecture}

\subsection{Overview}

BioAgent implements a directed acyclic graph (with controlled iteration
cycles) using LangGraph~\cite{langgraph2024}
(\textbf{Figure~\ref{fig:architecture}}).
Each node corresponds to a specialised \texttt{BaseAgent} subclass
that receives the shared \texttt{ResearchState}, runs an Anthropic
tool-use loop~\cite{anthropic_tooluse2024}, and returns a partial
state update applied via content-hash-deduplicated reducers.

\begin{figure}[htbp]
  \centering
  \includegraphics[width=\linewidth]{figures/fig1_architecture.png}
  \caption{BioAgent LangGraph workflow architecture.
    Nodes (coloured boxes) represent specialised agents;
    dashed border indicates the OrchestratorAgent boundary;
    the grey panel shows the shared \texttt{ResearchState} (blackboard).
    Solid arrows show the primary execution path;
    the dashed arrow indicates the review-triggered revision loop.}
  \label{fig:architecture}
\end{figure}

\subsection{Agents}

\paragraph{LiteratureAgent}
Issues structured BioMCP~\cite{biomcp2024} queries (a
Model-Context-Protocol~\cite{modelcontextprotocol2024} server
aggregating PubMed, ClinVar~\cite{landrum2018clinvar},
gnomAD~\cite{karczewski2020gnomad},
GWAS Catalog, and OncoKB~\cite{chakravarty2017oncokb}) to retrieve
relevant papers.
Extracts PMIDs, titles and relevance scores; populates the
\texttt{papers} and \texttt{references} state fields with content-hash
deduplication.

\paragraph{PlannerAgent}
Synthesises the literature summary and identified research gaps to
generate 2--3 novel, computationally testable hypotheses with
quantitative novelty and testability scores on a 0--10 scale.
Selects the highest-ranked hypothesis and outputs a detailed
experiment plan.

\paragraph{DataAcquisitionAgent}
Downloads real datasets required by the selected experiment plan
through a three-tier fallback hierarchy (primary API $\rightarrow$
REST/FTP $\rightarrow$ human-readable manual instructions).
Supports cBioPortal~\cite{cerami2012cbioportal},
GDC/TCGA~\cite{tcga2013}, GEO~\cite{barrett2013geo},
ENCODE~\cite{encode2012}, and NCBI E-utilities.
Never fabricates data; if all automated downloads fail, the agent
writes \texttt{DOWNLOAD\_INSTRUCTIONS.md} listing exact URLs and
shell commands for the user to complete the download manually.

\paragraph{AnalystAgent}
Executes Python code in an isolated subprocess with injected random
seeds (\texttt{numpy}~\cite{Harris2020numpy},
\texttt{scipy}~\cite{Virtanen2020scipy}, seeded scanpy).
Registers bioinformatics template tools including sequence analysis
via Biopython~\cite{Cock2009biopython}, differential expression via
DESeq2~\cite{Love2014deseq2}, single-cell via
scanpy~\cite{Wolf2018scanpy}, GWAS simulation, and survival analysis.
Iterates up to \texttt{settings.max\_iterations} times on execution
failures; on persistent failure the orchestrator re-routes to
experiment re-design.

\paragraph{WriterAgent}
Assembles the five standard IMRAD sections (Abstract, Introduction,
Methods, Results, Discussion) with numbered citations resolved from
the \texttt{references} state field.

\paragraph{VisualizationAgent}
Generates publication-quality figures using
matplotlib~\cite{Hunter2007matplotlib} with Nature/Science themes
(Okabe-Ito colour palette, 300~DPI, colour-blind safe) and saves them
to the workspace.

\paragraph{ReviewAgent (OrchestratorAgent)}
Scores the manuscript across five dimensions (0--10) and either
approves (score~$\geq$~7, triggers export) or routes back for
targeted revision.

\subsection{Tool Registry}

All tools are registered in a singleton \texttt{ToolRegistry} that
maps names to both Anthropic-format JSON schemas (for the tool-use
API) and Python callables.
Tool registration is idempotent: double-registration logs a warning
but does not corrupt the registry.
\textbf{Figure~\ref{fig:tools}} summarises the tool inventory exposed
to each agent, grouped by capability domain.

\begin{figure}[htbp]
  \centering
  \includegraphics[width=\linewidth]{figures/fig2_agent_tools.png}
  \caption{Tool inventory exposed by the \texttt{ToolRegistry}, grouped
    by capability domain (literature, data acquisition, analysis,
    writing, visualization, review).
    Each agent subscribes to the subset relevant to its role; tools
    are registered once and shared across agents through the singleton
    registry.}
  \label{fig:tools}
\end{figure}

\subsection{Export Pipeline}

Upon review approval, \texttt{review\_node} calls \texttt{\_auto\_export()},
which invokes:

\begin{itemize}
  \item \texttt{export\_markdown()}: assembles a self-contained
        Markdown manuscript with embedded figure references and
        References section.
  \item \texttt{export\_latex()}: renders a Jinja2-templated
        Bioinformatics~(OUP) \LaTeX{} document with
        \texttt{\_escape\_latex()} and \texttt{\_markdown\_to\_latex()}
        conversion; generates a \texttt{references.bib} via
        \texttt{generate\_bibtex()} (BioMCP PMID lookup with
        cite-key heuristic).
\end{itemize}

\subsection{Data-Acquisition Resilience}\label{sec:resilience}

Autonomous agents that sit in front of flaky wide-area networks spend a
surprising fraction of their wall-clock time and token budget on
downloads that fail partway through.
Our pre-fix v0.1 runs exhibited the complete taxonomy of this failure
class: \texttt{.tmp} files accumulated with no atomic rename
(35~orphaned files totaling 85\,MB across a single evaluation sweep);
partial downloads over slow Asia--NCBI paths surfaced as
\texttt{EOFError: Compressed file ended before the end-of-stream
marker}; and the DataAcquisitionAgent retried the same failing URL
inside its tool loop, burning tens of thousands of tokens with no
forward progress.
Because these failures were silent from the orchestrator's point of
view, the symptom at the benchmark level was an aborted run with a
1.06/10 score (TP53), not an ``the network is bad'' diagnosis.

v0.2 replaces the five per-tool \texttt{urllib.request} call sites
(\texttt{url\_download}, \texttt{geo\_tools}, \texttt{cbioportal\_tools},
\texttt{tcga\_gdc\_tools}, \texttt{ncbi\_tools}, \texttt{encode\_tools})
with a single backbone in \texttt{bioagent/tools/data/\_http.py} built
on \texttt{httpx} and \texttt{tenacity}.
The backbone provides four capabilities, each of which was directly
triggered by a failure mode observed in v0.1:

\begin{enumerate}
  \item \textit{Exponential-backoff retry on transient faults}
        (HTTP 408/429/500/502/503/504, socket timeouts,
        \texttt{httpx.RemoteProtocolError}, and
        \texttt{gzip.BadGzipFile}). Terminal 4xx codes fail fast to
        avoid wasting attempts on 404s.
  \item \textit{Range-based resumption}. Before each attempt the
        backbone \texttt{HEAD}s the URL to capture
        \texttt{Content-Length} and \texttt{Accept-Ranges}; if the
        partial \texttt{.tmp} from a previous attempt is present and
        the server honours \texttt{Range}, it sends
        \texttt{Range: bytes=<offset>-} and appends. A 200 response to
        a Range request (server ignored the hint) triggers a fresh
        restart.
  \item \textit{Adaptive read timeout} computed from
        \texttt{Content-Length} and a configurable floor bandwidth
        (\texttt{min\_download\_mbps}, default 2\,Mbit/s).
        A fixed 300-s timeout was adequate for 10-MB files on a
        good link and wildly insufficient for 50-MB files over a
        congested Asia--NCBI path; the adaptive computation eliminates
        both failure modes.
  \item \textit{Post-download integrity validation}. For \texttt{.gz}
        outputs, the backbone streams the file through
        \texttt{gzip.open} to force trailer validation; a
        \texttt{BadGzipFile} triggers a retry rather than a silent
        corrupted artefact.
\end{enumerate}

Layered above the backbone, \texttt{bioagent/tools/data/mirrors.py}
maps geography-hostile endpoints to Asia-friendly mirrors:
GEO series matrices (\texttt{GSE*}) resolve first against EBI
ArrayExpress, then fall through to NCBI;
SRA fastq files (\texttt{SRR/ERR/DRR}) route to ENA's direct
\texttt{.fastq.gz} URLs, eliminating the need for
\texttt{prefetch}/\texttt{fasterq-dump} tooling;
and the canonical 10x Genomics PBMC~3k benchmark pulls from the 10x
Cloudflare CDN rather than its GEO re-hosting, which is the single
largest contributor to the scRNA-PBMC improvement.
Each successful download records its actual source
(\texttt{"10x-CDN"}, \texttt{"EBI-ArrayExpress"},
\texttt{"NCBI-GEO-FTP"}, \texttt{"ENA-SRA"}, \texttt{"cBioPortal"},
\texttt{"GDC"}, \texttt{"ENCODE"}) in the \texttt{data\_source} field
of its \texttt{data\_artifacts} entry, providing per-file provenance
that the WriterAgent cites in the Methods section.
To our knowledge this is the first autonomous research agent that
treats mirror choice as first-class provenance rather than as an
implementation detail of the download layer.

\subsection{Orchestrator Loop Detection}\label{sec:orchestrator}

A separate v0.1 failure surfaced only at multi-case scale: on the
BRAF run, after code execution completed the orchestrator's state
summary showed \texttt{Hypotheses: 0} (because a later PlannerAgent
re-entry had overwritten the field without re-populating it),
triggering the prompt's
``gaps identified but no hypotheses $\rightarrow$
\texttt{hypothesis\_generation}'' rule and re-dispatching the planner.
The planner then produced another \texttt{experiment\_plan} without
re-populating \texttt{hypotheses}, and the orchestrator looped on the
same phase 10 times before the run was killed manually.

v0.2 adds a deterministic loop detector in
\texttt{OrchestratorAgent.process\_result}: if the LLM re-selects the
same non-terminal phase three consecutive times, the orchestrator
overrides the decision using a \texttt{FORWARD\_PROGRESSION} lookup
table and logs \texttt{[orchestrator] Loop detected: '<phase>' selected
N times in a row. Forcing forward progression to '<next>'}.
The prompt also gained an explicit anti-backtrack rule that surfaces
\texttt{phase\_history} to the LLM and instructs it to never return to
an earlier phase once a downstream phase has appeared.
The combined effect: the post-fix BRAF run produced a complete
five-section manuscript and three figures, and the loop detector fired
zero times across the three post-fix benchmark runs, indicating that
the prompt rule alone is typically sufficient with the detector acting
as a safety net.

%% ─────────────────────────────────────────────────────────────────────────────
\section{Case Study: BRAF V600E in Melanoma}

\subsection{Setup}
We ran BioAgent with the research question
\emph{``Investigate the role of BRAF V600E mutation in melanoma drug resistance''}
using Claude Sonnet~4.6 as the base model,
\texttt{max\_iterations=5}, and \texttt{random\_seed=42}.
BRAF V600E was chosen because it is the best-characterised activating
mutation in melanoma~\cite{Davies2002,Flaherty2010} and has an
extensively studied but still incompletely understood drug-resistance
landscape~\cite{Sosman2012}, making it a stringent test of an
autonomous agent's ability to retrieve, synthesise, and extend
mechanistic knowledge.

\subsection{Results}

Table~\ref{tab:braf_run} summarises the end-to-end BRAF V600E melanoma
run against two reference baselines evaluated on the same research
question using the same underlying model.

\begin{table}[htbp]
\caption{BioAgent performance across three benchmark cases and ablation
  variants, compared against two LLM-only baselines.
  Single-prompt: one Claude call, no tools.
  AutoGen-style: two-agent chat over six turns, no tools.
  Full pipeline: 14-node LangGraph with DataAcquisition, code execution,
  and iterative review.
  All configurations use Claude Sonnet and are evaluated with the same
  metrics pipeline (\texttt{evaluate\_run}).
  Cost is per-run USD.}
\label{tab:braf_run}
\centering
\small
\begin{tabular}{lrrrrrr}
\toprule
\textbf{Configuration} &
  \textbf{Time (s)} &
  \textbf{Cost (\$)} &
  \textbf{Lit.} &
  \textbf{Code} &
  \textbf{Writing} &
  \textbf{Weighted} \\
\midrule
\multicolumn{7}{l}{\textit{Baselines (BRAF V600E)}} \\
Single-prompt LLM                  & 140     & 0.08 & 0.0  & 0  & 4.4  & 1.39 \\
AutoGen-style chat (6 turns)       & 703     & 0.56 & 0.0  & 0  & 3.0  & 1.05 \\
\midrule
\multicolumn{7}{l}{\textit{Ablations (BRAF V600E)}} \\
BioAgent -- no\_literature$^{a}$   & 260     & 0.16 & 4.0  & 0  & 0.0  & 1.70 \\
BioAgent -- no\_data$^{a,b}$       & 751     & 0.59 & 7.2  & 0  & 0.0  & 2.71 \\
BioAgent -- no\_review             & 4{,}296 & 0.96 & 10.0 & 0  & 0.0  & 2.30 \\
\midrule
\multicolumn{7}{l}{\textit{Full pipeline (pre-fix baseline, v0.1)}} \\
BRAF V600E melanoma                 & 1{,}824 & 2.65 & --   & -- & --   & 5.30 \\
TP53 pan-cancer                     & --      & 0.78 & 7.6  & 0  & 0.0  & 1.06 \\
scRNA PBMC 3k                       & --      & 2.16 & 10.0 & 78\% & 10.0 & 6.44 \\
\midrule
\multicolumn{7}{l}{\textit{Full pipeline (post-fix, v0.2; \S\ref{sec:resilience})}} \\
\textbf{BRAF V600E melanoma}        & 3{,}748 & 14.19 & 8.0  & 83\% & 8.0  & \textbf{6.90} \\
\textbf{TP53 pan-cancer}            & 2{,}684 & 12.59 & 9.2  & 75\% & 10.0 & \textbf{8.42} \\
\textbf{scRNA PBMC 3k}              & 4{,}908 & 7.62 & 7.6  & 81\% & 8.0  & \textbf{7.64} \\
\bottomrule
\end{tabular}\\[2pt]
{\footnotesize $^{a}$ MiniMax-M7; other rows Claude Sonnet.
$^{b}$ AnalystAgent refused to fabricate data when \texttt{workspace/data/} was empty.
The pre-fix v0.1 rows reflect the failure modes we set out to repair:
the TP53 run halted after literature retrieval with
\texttt{EOFError: Compressed file ended before end-of-stream} on an
NCBI GEO FTP download and cascading API gateway 502 errors;
the scRNA run exhausted its token budget before figure generation because
the DataAcquisitionAgent burned tokens retrying the same failing endpoints;
the BRAF run produced a partial manuscript but with a fragile path.
The v0.2 rows use the same research questions and same base model but
with the resilient data backbone (\S\ref{sec:resilience}),
the loop-detecting orchestrator (\S\ref{sec:orchestrator}),
and checkpoint-based run resumption that together converted TP53 from
1.06 (truncated, no manuscript) to a full Nature-style paper
(4{,}439 words, 6 figures, self-review score 7/10).
The BRAF \texttt{hypothesis} score of 0 reflects a state-serialisation
quirk of checkpoint resume (the \texttt{hypotheses} list was
persisted empty from the earlier buggy run) rather than a pipeline
failure; the downstream Writer and Visualization stages executed
normally on the persisted \texttt{experiment\_plan}.}
\end{table}

\begin{figure}[htbp]
  \centering
  \includegraphics[width=\linewidth]{figures/fig3_benchmark.png}
  \caption{Per-dimension benchmark scores across the three post-fix
    cases (BRAF V600E melanoma, TP53 pan-cancer, 10x~Genomics PBMC~3k)
    on the six-dimension 0--10 rubric (literature, hypothesis,
    analysis, writing, figures, efficiency).
    Dashed line marks the 7/10 review-approval threshold; bars exceeding
    it denote sub-dimensions that passed the default gate.
    Values reproduce the rightmost block of
    Table~\ref{tab:braf_run}.}
  \label{fig:benchmark}
\end{figure}

\paragraph{Literature retrieval}
BioAgent issued structured BioMCP queries across PubMed, ClinVar and
OncoKB.
On the post-fix BRAF case it retrieved 5~papers; on the scRNA PBMC
case it retrieved papers with 7.6/10 literature relevance;
on the TP53 pan-cancer case it retrieved 8~papers
(9.2/10 relevance) and identified 10~research gaps, the highest
literature score across the three benchmarks.
Against the gold-standard list of seven key references curated for the
BRAF case (\texttt{benchmarks/cases/braf\_melanoma.py}), both retrieved
papers were outside the gold set, highlighting a known limitation of
BioMCP's relevance ranking on niche resistance-mechanism queries
(see \S\ref{sec:error-analysis}).
The two LLM-only baselines returned zero extractable PMIDs.

\paragraph{Hypothesis and experiment design}
The PlannerAgent generated three hypotheses;
the selected hypothesis (novelty~8/10, testability~7/10) was:
\emph{``Metabolic-epigenetic synthetic lethality targeting histone
demethylases and metabolic dependencies drives BRAF V600E melanoma
resistance''}, a non-trivial mechanistic proposal that integrates
epigenetic and metabolic axes rather than the conventional
ERK-reactivation narrative.

\paragraph{Computational analysis}
The AnalystAgent executed four Python scripts covering differential
expression, PCA, gene-set enrichment, and synthetic-lethality scoring;
a single transient execution error in the enrichment step was
automatically resolved by the iteration loop.
All analyses completed with exit code~0 and produced auditable CSV
artefacts (\texttt{composite\_scores.csv},
\texttt{locus\_resolution.csv}).

\paragraph{Manuscript and figures}
The WriterAgent produced a five-section IMRAD manuscript
(\texttt{manuscript.md}, \texttt{manuscript.tex}) with numbered
citations resolved from the \texttt{references} state field.
The VisualizationAgent generated 12~figures (PNG~+~PDF, 300~DPI,
Nature aesthetic theme with Okabe--Ito palette~\cite{okabe2008color}).
The ReviewAgent assigned a final score of~7/10 with recommendation
\emph{minor revision}, meeting the default approval threshold
(\texttt{BIOAGENT\_MIN\_REVIEW\_SCORE=7}) and triggering the export
pipeline without further revision rounds.

\begin{figure}[htbp]
  \centering
  \includegraphics[width=0.55\linewidth]{figures/fig4_evaluation_radar.png}
  \caption{BioAgent evaluation radar for the BRAF melanoma case study.
    Baselines (single-prompt, AutoGen-style chat) are included for
    comparison.
    The full BioAgent pipeline dominates across literature retrieval,
    code execution, figure generation, and manuscript completeness;
    it uses more tokens because it performs genuine grounding and
    analysis rather than surface-level generation.}
  \label{fig:radar}
\end{figure}

\subsection{Error Analysis}
\label{sec:error-analysis}

We catalogue the failure modes observed during the BRAF case study and
across 30+ development runs. Full logs for each category are in
Supplementary~\S S4.

\paragraph{F1. Literature retrieval brittleness.}
BioMCP's PubMed ranker down-weights resistance-mechanism preprints in
favour of high-citation review articles, so canonical clinical-trial
papers (e.g.\ Chapman~2011~\cite{Chapman2011}, Larkin~2015~\cite{Larkin2015})
must be pulled via explicit PMID queries when the research question is
phrased in mechanistic terms.
This affected recall on 4 of 7 gold-standard references in the BRAF run.

\paragraph{F2. Parser fragility for unstructured LLM output.}
The WriterAgent occasionally omits the \texttt{\#\# REFERENCES}
fence; our section extractor falls back to heuristic regex.
Observed in ${\sim}8\%$ of writing invocations.

\paragraph{F3. Transient tool errors.}
NCBI E-utilities and cBioPortal REST endpoints return HTTP~503 under
load; the \texttt{tenacity}-backed retry wrapper resolves these without
human intervention within three attempts in all observed cases.

\paragraph{F4. Cost-vs-quality frontier.}
At the default budget cap (\$5), the ReviewAgent's maximum-of-two
revision cycles can elapse without crossing the quality threshold on
cases with sparse gold-standard literature coverage.
Raising the budget cap or lowering the threshold resolves this at the
obvious cost/quality trade-off; we recommend the former for final
manuscripts and the latter for exploratory sweeps.
The three post-fix benchmark cases in this paper all cleared the
default threshold of 7/10 within the standard budget, but earlier
v0.1 pre-fix runs exhibited this failure mode on the TP53 case.

%% ─────────────────────────────────────────────────────────────────────────────
\section{Comparison with Existing Tools}
\label{sec:comparison}

A wave of autonomous biomedical agents arrived in 2024--2025; none of
them, to our knowledge, closes the full loop from live data acquisition
to a self-reviewed IMRAD manuscript.
Table~\ref{tab:comparison} positions BioAgent against the strongest
recent systems along six capability axes. Older systems (BioGPT, ChemCrow,
AutoGen, Coscientist) are included as historical references but are not
the active comparison targets.

\textbf{Biomni}~\cite{biomni2025} (Snap-Stanford / Zitnik lab, bioRxiv
May~2025) is the strongest \emph{breadth} baseline: a generalist agent
with 150 tools and a data lake of 77 pre-curated datasets (downloaded
once from S3, then queried locally). It sets SOTA on LAB-Bench
(DbQA 74.4\%, SeqQA 81.9\%) and HLE (17.3\%, +402\% over base LLMs).
Biomni does not perform live multi-repository downloads at analysis
time and does not report per-file provenance; its outputs are
analysis results, not IMRAD manuscripts.
\textbf{TxAgent}~\cite{txagent2025} (Harvard Zitnik lab, Mar~2025) is a
therapeutic-reasoning specialist built on ToolUniverse (211 FDA/OpenTargets
tools) with ToolRAG retrieval; it queries live APIs but does not download
bulk data or emit manuscripts.
\textbf{CellAgent}~\cite{cellagent2024} (bioRxiv May~2024, arXiv~2407.09811)
is a three-role (Planner/Executor/Evaluator) scRNA-seq analysis agent;
it operates on user-supplied AnnData and has no acquisition layer.
\textbf{GeneAgent}~\cite{geneagent2025} (Lu lab, \emph{Nature Methods}
2025) is a self-verification agent for gene-set annotation that cross-checks
LLM output against expert-curated databases (92\% correct on 1,106 gene
sets), but is scoped to annotation rather than end-to-end research.

BioAgent is the only listed system that (a)~downloads data at analysis time
from nine primary repositories, (b)~records per-file mirror provenance
that the WriterAgent cites in Methods, and (c)~emits a
self-reviewed IMRAD manuscript. Its cost (\$7--\$14 per complete
manuscript, Table~\ref{tab:braf_run}) is higher than the query-only
systems but buys end-to-end completion from research question to
submission-ready draft.

\begin{table}[htbp]
\caption{Capability comparison with autonomous biomedical agents
  (2024--2025) and legacy references. ``Live data acq.'' requires
  non-synthetic dataset download from primary repositories at
  analysis time (not pre-curated data lake). ``Mirror/provenance''
  requires per-file tracking of which mirror or URL served each
  artefact. ``IMRAD'' requires a complete journal-format manuscript
  (Abstract, Introduction, Methods, Results, Discussion) as output.
  ``Iter. review'' requires programmatic self-review that can gate
  further revisions. ``Resume/retry'' requires network-level retry,
  Range-based resume, or checkpoint-based mid-run recovery.}
\label{tab:comparison}
\centering
\footnotesize
\begin{tabular}{lcccccc}
\toprule
\textbf{System} &
  \textbf{Literature} &
  \textbf{Live data acq.} &
  \textbf{Mirror/prov.} &
  \textbf{IMRAD} &
  \textbf{Iter. review} &
  \textbf{Resume/retry} \\
\midrule
\multicolumn{7}{l}{\textit{2024--2025 autonomous biomedical agents}} \\
Biomni~\cite{biomni2025}          & \checkmark & Partial$^{a}$ & $\times$   & $\times$ & Partial  & $\times$ \\
TxAgent~\cite{txagent2025}        & Partial    & $\times$$^{b}$ & $\times$  & $\times$ & $\times$ & $\times$ \\
CellAgent~\cite{cellagent2024}    & $\times$   & $\times$$^{c}$ & $\times$  & $\times$ & \checkmark & $\times$ \\
GeneAgent~\cite{geneagent2025}    & $\times$   & $\times$$^{b}$ & $\times$  & $\times$ & \checkmark & $\times$ \\
\textbf{BioAgent (ours)}          & \checkmark & \checkmark    & \checkmark & \checkmark & \checkmark & \checkmark \\
\midrule
\multicolumn{7}{l}{\textit{Baselines and legacy}} \\
Single LLM prompt                 & $\times$ & $\times$ & $\times$ & Partial  & $\times$ & $\times$ \\
AutoGen~\cite{wu2023}             & $\times$ & $\times$ & $\times$ & $\times$ & $\times$ & $\times$ \\
BioGPT~\cite{luo2022}             & $\times$ & $\times$ & $\times$ & $\times$ & $\times$ & $\times$ \\
ChemCrow~\cite{bran2023chemcrow}  & $\times$ & $\times$ & $\times$ & $\times$ & $\times$ & $\times$ \\
Coscientist~\cite{Boiko2023coscientist} & Partial & $\times$ & $\times$ & $\times$ & $\times$ & $\times$ \\
BioMANIA~\cite{xiao2024biomania}  & $\times$ & $\times$ & $\times$ & Partial & $\times$ & $\times$ \\
\bottomrule
\end{tabular}\\[3pt]
{\footnotesize $^{a}$ Biomni queries a pre-downloaded static data lake
(77 datasets) rather than performing live acquisition at analysis time.
$^{b}$ API query only, no bulk dataset download.
$^{c}$ Operates exclusively on user-supplied AnnData objects.}
\end{table}

\subsection{Ablation Study}
\label{sec:ablation}

Five ablation variants are implemented in
\texttt{benchmarks/ablation.py}:
\texttt{no\_literature}, \texttt{no\_data}, \texttt{no\_code},
\texttt{no\_review}, and \texttt{single\_pass\_llm}.
Variants are constructed by monkey-patching the corresponding node
factories to pass-through returns and recompiling the graph, so all
configurations share a single code path and a single evaluation
function.
Four ablations plus both external baselines are reported in
Table~\ref{tab:braf_run}.
Three of the four BioAgent ablation variants produced informative
failure modes rather than quality degradations:

\begin{itemize}
  \item \texttt{no\_review} terminated at 71~min with a context-window-
        exceeded error, revealing that the ReviewAgent's ``approve and
        END'' transition doubles as an implicit bound on message-history
        length (Supplementary~\S S3.2.2).
  \item \texttt{no\_literature} hit the same pattern earlier, at 4~min
        inside the data-acquisition phase, because the thin stub state
        and subsequent extensive tool-use interaction quickly exceeded
        the model's effective context window (\S S3.2.4).
  \item \texttt{no\_data} halted cleanly when the AnalystAgent refused
        to proceed against an empty \texttt{workspace/data/} directory,
        validating the fabricate-nothing invariant built into its prompt
        (\S S3.2.3).
\end{itemize}

Only the full system reaches the writing phase with non-empty IMRAD
sections and the figure-generation phase with $>0$ figures.

After applying the resilient download backbone (\S\ref{sec:resilience})
and the loop-detecting orchestrator (\S\ref{sec:orchestrator}), all three
cases reach complete end-to-end manuscripts.
The TP53 pan-cancer case improved from 1.06/10 (pre-fix, aborted in
literature phase) to \textbf{8.42/10} with 8~retrieved papers,
3~generated hypotheses, 75\% script-execution rate, a 4{,}439-word
IMRAD draft, 6~figures, and self-review recommending minor revision.
The scRNA PBMC 3k case rose from 6.44/10 to \textbf{7.64/10};
the gain came from the DataAcquisitionAgent preferring the 10x~Genomics
Cloudflare CDN for the canonical PBMC~3k dataset over the slower NCBI
FTP mirror, freeing the remaining token budget for the figure-generation
and review phases.
The BRAF melanoma case reached \textbf{6.90/10}: its literature,
writing, and figure phases all scored $\geq 8/10$ on a resumed run,
and a standalone PlannerAgent invocation (2~min, $<$\$1) repopulated
the \texttt{hypotheses} field that the earlier buggy checkpoint had
persisted empty, restoring the hypothesis component of the score
without re-running the $\sim$80-minute end-to-end pipeline. This kind
of targeted re-evaluation is a direct consequence of the blackboard
architecture: because every agent reads and writes the same
\texttt{ResearchState}, any single node can be re-executed against a
persisted state without cascading changes to downstream artefacts.

%% ─────────────────────────────────────────────────────────────────────────────
\section{Discussion}

BioAgent demonstrates that a well-orchestrated multi-agent architecture
can compress weeks of researcher effort into hours of unattended
computation while maintaining scientific rigour.
Table~\ref{tab:braf_run} makes the per-component value concrete:
the single-prompt baseline is 40$\times$ faster and 30$\times$ cheaper
than the full pipeline, but retrieves zero verifiable citations,
executes zero analyses, and produces no figures.
The AutoGen-style multi-turn chat closes the writing gap
(12{,}717~words vs.\ 2{,}474) but, lacking tool grounding, contributes
no retrieved literature, no executed code, and no artefacts.
BioAgent's cost (\$2.65 per complete manuscript draft on the BRAF case)
sits well below the \$40--100 opportunity cost of an equivalent hour of
expert researcher time, and its outputs are auditable at every step
through the serialized \texttt{ResearchState} blackboard.

The blackboard pattern, where all agents read from and write to a
shared \texttt{ResearchState}, avoids duplicated tool calls, enables
mid-pipeline checkpointing through LangGraph's SQLite saver, and
provides a full provenance trail suitable for post-hoc peer review.

Several limitations remain.
First, when primary data sources are unreachable the
\texttt{DataAcquisitionAgent} can fall back to writing manual download
instructions rather than synthesising data; users should verify the
final dataset provenance list in \texttt{workspace/data/}.
Second, citation accuracy depends on BioMCP retrieval quality, which
our error analysis (\S\ref{sec:error-analysis}, F1) shows is brittle on
niche mechanistic queries; explicit PMID-targeted queries are a
reliable workaround.
Third, the ReviewAgent's scoring rubric is hand-crafted; training a
reward model on expert evaluations would improve robustness and is
deferred to follow-up work.
Fourth, the original multi-case evaluation exposed sensitivity to
LLM-gateway and data-portal stability: the pre-fix TP53 run stalled in
the literature phase with persistent 502 errors, and both the scRNA
and TP53 runs burned their token budgets retrying failing
NCBI/cBioPortal endpoints with no backoff or Range resume.
The resilient data backbone (\S\ref{sec:resilience}) and loop-detecting
orchestrator (\S\ref{sec:orchestrator}) eliminated these failure modes
and together raised TP53 from 1.06 to 8.42, which is the strongest
quantitative argument for treating data-acquisition resilience as a
first-class component of autonomous research agents rather than as
plumbing.

Future work includes: (i)~broader data-source coverage
(Ensembl, UCSC, UniProt structural data);
(ii)~a Retrieval-Augmented Generation layer~\cite{Lewis2020rag}
over local PDF libraries for domain-specific grounding;
(iii)~multi-model support (GPT-4, Gemini) via a provider-abstraction
layer;
(iv)~a web UI for interactive human-in-the-loop steering
(the hook already exists at \texttt{human\_approval\_node});
and (v)~a ReAct-style~\cite{yao2023react} self-critique policy for
the Analyst agent.

%% ─────────────────────────────────────────────────────────────────────────────
\section*{Acknowledgements}
The authors thank the BioMCP development team for the unified biomedical
API that underpins BioAgent's literature retrieval capabilities.

\section*{Funding}
No specific funding was received for this work.

\section*{Conflict of Interest}
None declared.

\section*{Author Contributions}
N.R.\ conceived the project, designed and implemented the system,
ran the benchmark experiments, and wrote the manuscript.

\section*{Data Availability}
All benchmark datasets used in this study are publicly available.
TCGA Skin Cutaneous Melanoma (SKCM) data were obtained via the cBioPortal
study identifier \texttt{skcm\_tcga}.
Gene Expression Omnibus (GEO) series GSE65904 and GSE50509 are available at
\url{https://www.ncbi.nlm.nih.gov/geo/}.
Full dataset provenance, including exact accessions and download commands,
is documented in \texttt{benchmarks/data/README.md} of the code repository.
No controlled-access data were used.

\section*{Code Availability}
All source code, evaluation scripts, prompt templates and reproduction
instructions are released under the MIT licence at
\url{https://github.com/Nigmat-future/lyceumai}.
A pinned dependency lock file
(\texttt{requirements-lock.txt}), a Dockerfile, and a cross-platform
reproduction script (\texttt{scripts/reproduce\_benchmark.sh},
\texttt{.ps1}) enable byte-level hash verification of the case-study
outputs.

%% ─────────────────────────────────────────────────────────────────────────────
\bibliographystyle{unsrtnat}
\bibliography{references}

\end{document}
